\section{Conclusion}

This work establishes a fundamental connection between neural architecture design and the effectiveness of biological attention mechanisms. We introduce the \textit{compression-covariance dichotomy}: sequential architectures (RNN/LSTM) compress information through a bottleneck that limits biological enhancement, while transformers preserve pairwise relationships that biological mechanisms can directly modulate.

Our experiments confirm this hypothesis quantitatively. Bio-enhanced transformers improve by 11.7\% over baselines, compared to 6.2\% for LSTM and 4.2\% for RNN. This ordering---Transformer > LSTM > RNN---reflects each architecture's capacity to preserve the relational structure that biological mechanisms leverage.

Key contributions of this work include:
\begin{itemize}
\item A theoretical framework (compression vs. covariance) explaining differential effectiveness of biological attention across architectures
\item Implementations of four biologically-inspired mechanisms (biased competition, lateral inhibition, tuning sharpening, noise decorrelation) for RNN, LSTM, and Transformer
\item Empirical demonstration that transformers benefit most from biological enhancement, with ablation studies quantifying individual mechanism contributions
\item Analysis revealing emergent sparsity, improved gradient flow, and increased multi-head diversity in bio-enhanced transformers
\end{itemize}

These findings have practical implications for neural architecture design: preserving relational structure enables more effective integration of biological computational principles. They also suggest a research direction: developing architectures that combine the efficiency of recurrence with the relational preservation of attention, potentially achieving the best of both paradigms.

Looking forward, we anticipate that biological attention mechanisms will become increasingly important as models scale and efficiency becomes paramount. The emergent sparsity we observe---arising naturally from biologically-motivated computations---offers a path toward more efficient attention without sacrificing flexibility. As the field continues to draw inspiration from neuroscience, understanding \textit{why} certain architectures benefit more from biological principles will be essential for informed design choices.
